% SPDX-License-Identifier: Apache-2.0
\section{Comparison}
\label{sec:comparison}

Table~\ref{tab:compare} compares the approaches. ``Crostini'' means that the
approach works in Crostini with no root on ChromeOS, no kernel module, and no
developer mode. ``Both'' means that the key works on the Chromebook and on the
remote host at the same time.

\begin{table*}[t]
\centering
\caption{Comparison of the approaches. Y: yes. N: no. P: partly. ``a'' marks an assumption.}
\label{tab:compare}
\scriptsize
\setlength{\tabcolsep}{3pt}
\begin{tabular}{@{}p{24mm}p{30mm}ccccccp{38mm}p{30mm}@{}}
\toprule
Approach & Existing software & FIDO2 & OTP & PIV & PGP & Crostini & Both & Main security risk & Timeout risk \\
\midrule
A. USB/IP from Crostini & jiegec/usbip (serial, bug \#51); usbipdcpp (async) & Y\,a & Y\,a & Y\,a & Y\,a & P & N & Remote root has raw USB to all applets: reset, PIN exhaustion, OTP changes. No transport authentication. & High for jiegec/usbip; low for usbipdcpp (a) \\
B. CTAPHID relay, uhid & fido-forward (remote side as is); new usbfs key side & Y & N & N & N & P & N & Remote asks for assertions for any site. Reset and credential management blocked by the new policy. & Low: keepalives pass \\
B$'$. qubes-app-u2f over SSH & qubes-app-u2f & P & N & N & N & N & N & Built-in command filter & High: no keepalives \\
C. Virtual device, delegated crypto & virtual-fido, tpm-fido & N & N & N & N & -- & -- & Not possible: a real key cannot sign for it & -- \\
D. Agent forwarding & OpenSSH, GnuPG, p11-kit & N & N & Y & Y & P & N & Remote uses the agent while the session is open & Low \\
E1. Hybrid (phone) & Chrome & phone & N & N & N & -- & Y & Low & Needs Bluetooth at the remote host \\
E2. CRD forwarding & CRD & internal & N & N & N & -- & -- & -- & Not available \\
E3. WebAuthn proxy extension & auth-relay & Y\,a & N & N & N & N & N & No origin in the proxied request; origin is guessed & Low \\
E4. Smart card relay & vsmartcard & N & N & Y\,a & Y\,a & P & N & Remote sends any APDU & Low (a) \\
E5. Chromebook Chrome via SOCKS & Chrome, OpenSSH & Y & N & N & N & Y & Y & Low: the key stays local & Low; only for web apps the remote serves \\
E6. \code{vmc key-attach} & ChromeOS & U2F & N & N & N & P & Y & As B & No FIDO2 \\
\bottomrule
\end{tabular}
\end{table*}

Three approaches drop out at once. C cannot work, because a real key does not
sign for another authenticator. E1 and E2 are not available to this user. E6
lacks FIDO2. Of the rest, only A and B serve goal~1 inside the remote
desktop, and both depend on Findings~\ref{find:nohidraw}
and~\ref{find:detach}.

A forwards everything, and B forwards only FIDO. That is the main difference
in risk. Under A, the remote root user can reset the PIV and OpenPGP applets,
use up their PIN tries, and rewrite the OTP slots. Under B, the remote root
user can reach only the FIDO interface, and the relay can refuse the commands
that change the key.
